\documentclass[12pt,a4paper]{article}

\usepackage[T2A]{fontenc}
\usepackage[utf8]{inputenc}
\usepackage[russian]{babel}
\usepackage{microtype}
\usepackage{amsmath,amssymb}
\usepackage[margin=2.5cm]{geometry}
\usepackage{parskip}
\usepackage{titlesec}
\usepackage{enumitem}
\usepackage{xcolor}

\definecolor{slidegray}{RGB}{80,80,80}

\titleformat{\section}{\large\bfseries}{}{0em}{}[\vspace{-4pt}\rule{\linewidth}{0.4pt}]
\titlespacing{\section}{0pt}{16pt}{6pt}

\setlist[itemize]{itemsep=3pt,topsep=4pt}

\begin{document}

\begin{center}
  {\LARGE\bfseries Текст для докладчика}\\[6pt]
  {\large Могут ли простые СДУ описывать временные ряды фондового рынка?}\\[4pt]
  {\color{slidegray} Дмитрий Болт\quad---\quad Май 2026}
\end{center}

\bigskip

% ================================================================
\section{Слайд 1. Титульный}
% ================================================================

Это учебный мини-исследовательский проект.
Моя цель --- не построить производственную торговую модель.
Я использую эталонные стохастические дифференциальные уравнения
из учебника Оксендаля как математический язык
и задаю вопрос: какой из этих эталонов описывает
какое именно преобразование реальных данных о ценах акций, и в каком смысле?

% ================================================================
\section{Слайд 2. Постановка задачи и финансовая постановка}
% ================================================================

Хочу убедиться, что финансовая постановка понятна, поскольку мы работаем
с реальными данными.

\textbf{Акция} --- это торгуемый финансовый инструмент, доля собственности
в компании.
Её \textbf{цена закрытия} ($P_t^k$) --- это последняя цена сделки
в конце каждого торгового дня.
Рынок открыт примерно \textbf{252 дня в году} ---
выходные и праздники исключены.

Ключевой методологический момент: один и тот же ценовой ряд порождает
пять структурно различных эмпирических объектов.
Модель может хорошо работать для одного преобразования и плохо --- для другого.
Это организующий принцип всего исследования.

% ================================================================
\section{Слайд 3. Три учебных эталона}
% ================================================================

Все три эталона взяты непосредственно из Оксендаля.

\textbf{Диффузия Ито} (определение 4.1.1, уравнение 4.1.6) --- наиболее общий случай.
Она говорит: процесс изменяется за счёт детерминированного члена сноса
и случайного диффузионного члена.
Для дневных лог-доходностей это естественный язык:
доходность каждого дня --- это малый зашумлённый шаг.

\textbf{Геометрическое броуновское движение}
($\text{GBM}$, глава 5.1, уравнение 5.1.5)
специализирует снос пропорционально текущему уровню,
что даёт мультипликативный рост.
В финансах ($N_t$) --- это уровень цены акции или портфеля.
Решение --- экспонента с броуновским компонентом:
стандартная картина случайного блуждания со сносом для цен.

\textbf{Процесс Орнштейна--Уленбека}
($\text{OU}$, упражнение 5.7)
вводит возврат к среднему:
если процесс отклоняется выше долгосрочного уровня~$m$,
член сноса ($\theta(m - X_t)$) возвращает его обратно.
Здесь скорость возврата к среднему --- ($\theta$).
Именно такое поведение мы ожидаем от детрендированного ряда,
осциллирующего вокруг долгосрочного среднего.

% ================================================================
\section{Слайд 4. Данные и конвейер обработки}
% ================================================================

У нас 94 акции, дневные наблюдения,
примерно 820 точек данных на акцию.
Три представительных акции ---
NVDA (Nvidia), META (Meta Platforms), GOOG (Alphabet/Google) ---
это крупные американские технологические компании,
отобранные по явным критериям скрининга, показанным на следующем слайде.

\textbf{Шаг детрендирования} --- ключевая операция:
мы подгоняем простой линейный тренд методом наименьших квадратов ($\text{OLS}$)
к каждому ряду лог-цен и оставляем остаток.
Этот остаток ($\widetilde{X}_t^k$) ---
это объект, который мы проверяем на соответствие эталону Орнштейна--Уленбека.
Без детрендирования цены акций со временем растут,
и эталон ОУ очевидно не применим.

% ================================================================
\section{Слайд 5. Представительные акции: скрининг}
% ================================================================

Я выбрал три акции по трём критериям.

\textbf{Коэффициент детерминации линейного тренда} ($R^2$):
насколько хорошо линейный тренд объясняет сырую лог-цену?
Высокое значение означает, что ряд определяется сносом ---
хорошо для эталона геометрического броуновского движения.

\textbf{Автокорреляция при лаге один} ($\text{ACF}$ при лаге 1)
детрендированного уровня: насколько персистентен детрендированный ряд ($\widetilde{X}_t^k$)?
Высокое значение означает медленный возврат к среднему ---
хорошо для эталона ОУ.

\textbf{Стандартное отклонение детрендированного ряда} ($\sigma$):
амплитуда детрендированных осцилляций ---
чем больше, тем нагляднее визуальная картина.

NVDA имеет наиболее сильный тренд и наибольшую амплитуду.
META показывает наибольшую персистентность детрендированных уровней.
GOOG спокойнее, но по-прежнему хорошо выражен.
AAPL слабее по всем метрикам и служит резервным вариантом.

% ================================================================
\section{Слайд 6. Сырые лог-цены vs детрендированные лог-цены}
% ================================================================

Это ключевой визуальный образ исследования.

\textbf{Левые панели} показывают сырые лог-цены:
они растут со временем, иногда резко,
и выглядят в точности как траектории геометрического броуновского движения ---
снос-доминированные с аддитивным шумом поверх.

\textbf{Правые панели} показывают детрендированные остатки.
После удаления линейного тренда ряд осциллирует вокруг нуля.
Это визуально гораздо более совместимо с эталоном Орнштейна--Уленбека.

Ключевой вывод: одни и те же данные дают обе картины в зависимости от преобразования.

% ================================================================
\section{Слайд 7. Распределения приращений: лог-доходности}
% ================================================================

\textbf{Лог-доходность} ($\Delta X_{t_i}^k = \ln P_{t_i}^k - \ln P_{t_{i-1}}^k$) ---
это стандартная мера дневного изменения цены в финансах.
Она симметрична и приблизительно масштабно-инвариантна:
удвоение цены даёт одинаковую доходность независимо от начального уровня.

Гистограммы показывают, что детрендирование незначительно смещает центр
распределения к нулю, но не меняет форму шума.
Форма шума --- это диффузионный компонент;
смещение среднего --- это компонент сноса.
Язык диффузии Ито описывает оба.

% ================================================================
\section{Слайд 8. Скользящие диагностики: постоянна ли диффузия?}
% ================================================================

Эталон Ито (уравнение 4.1.6) с постоянным сносом $u$ и постоянным коэффициентом
диффузии $v$ предсказывает, что дневные доходности имеют постоянное среднее
и постоянное стандартное отклонение ($\sigma$) во времени.

21-дневное скользящее среднее (левая панель) остаётся малым и флуктуирует вблизи нуля ---
снос трудно обнаружить в дневных данных, что согласуется с зашумлённой диффузией.

21-дневное скользящее стандартное отклонение (правая панель) явно меняется со временем:
\begin{itemize}
  \item NVDA: отношение максимальной к минимальной скользящей волатильности
        составляет примерно 3.7 раза.
  \item META: отношение примерно 6.8 раза.
  \item GOOG: отношение примерно 3.1 раза.
\end{itemize}

Это хорошо известный эмпирический факт в финансах, называемый
\textbf{кластеризацией волатильности} ---
рынок чередует спокойные и турбулентные режимы.
Стохастическое дифференциальное уравнение с постоянными коэффициентами
не может этого уловить.
Оно остаётся полезным эталоном, но имеет известное, измеримое ограничение.

% ================================================================
\section{Слайд 9. Персистентность ОУ: АКФ детрендированных уровней}
% ================================================================

Теория Орнштейна--Уленбека (упражнение 5.7) предсказывает, что
автокорреляция детрендированного ряда ($\widetilde{X}_t$) убывает как ($e^{-\theta\ell}$),
где скорость возврата к среднему равна ($\theta$),
а лаг --- ($\ell$).

Эмпирические графики функции автокорреляции показывают именно этот паттерн:
высокие положительные значения при лаге один ---
0.9916 для NVDA, 0.9919 для META, 0.9880 для GOOG ---
которые убывают постепенно, а не обнуляются немедленно.

Это не доказывает точный закон ОУ.
Но это эмпирический признак, мотивирующий приближение ОУ.
Медленное геометрическое убывание согласуется со значением ($\hat\theta$)
около 0.012, что соответствует периоду полуравновесия примерно 57 торговых дней.

% ================================================================
\section{Слайд 10. Детрендированные уровни и скользящее среднее}
% ================================================================

Этот слайд показывает медленный, волнообразный характер возврата к среднему.

63-дневное скользящее среднее (синяя линия) показывает, что локальный
долгосрочный уровень детрендированного ряда ($\widetilde{X}_t^k$)
сам медленно меняется со временем.
Это согласуется с очень медленным возвратом к среднему:
процесс не возвращается резко к фиксированному равновесию,
а плавно дрейфует к медленно эволюционирующему локальному среднему.

% ================================================================
\section{Слайд 11. Оценка параметров ОУ: прокси AR(1)}
% ================================================================

Мост между непрерывным стохастическим дифференциальным уравнением ОУ
и дискретными дневными данными --- это прокси авторегрессионной модели первого порядка (AR(1)).

Если процесс удовлетворяет уравнению ОУ, то точное дискретное решение
за один день имеет вид:
($\widetilde{X}_{t+1} = m(1-e^{-\theta}) + e^{-\theta}\widetilde{X}_t + \varepsilon_t$).

Мы оцениваем это как ($\hat\alpha + \hat\phi\,\widetilde{X}_t$).
Из подогнанного ($\hat\phi$)
мы восстанавливаем ($\hat\theta$) как минус логарифм от ($\hat\phi$),
то есть ($\hat\theta = -\ln\hat\phi$).

\textbf{Период полуравновесия} ($t_{1/2} = \ln 2\,/\,\hat\theta$) ---
это время, за которое отклонение от среднего уменьшается вдвое.
Для всех трёх акций период полуравновесия составляет примерно 57 торговых дней ---
то есть примерно \textbf{три календарных месяца}.
Это очень медленный возврат к среднему:
не резкий отскок, а плавный многомесячный дрейф к долгосрочному уровню.

% ================================================================
\section{Слайд 12. Симуляции: связь теории и данных}
% ================================================================

Этот слайд замыкает петлю между теоретическими эталонами и эмпирическими данными.

\textbf{Слева:} чистое броуновское движение --- симметричное случайное блуждание,
без сноса, без возврата к среднему.

\textbf{По центру:} геометрическое броуновское движение (GBM) ---
тот же броуновский шум, умноженный на текущий уровень,
что даёт экспоненциальный рост при случайных возмущениях.
Это визуальная модель для сырых цен акций.

\textbf{Справа:} процесс Орнштейна--Уленбека, откалиброванный на GOOG ---
($\hat\theta = 0.0121$),
($\hat m = 0.0207$),
($\hat\sigma = 0.0187$).
Заметьте, как траектории осциллируют вокруг пунктирного долгосрочного среднего,
не уходя в бесконечность.
Возврат медленный и плавный --- согласующийся с периодом полуравновесия 57 дней.

% ================================================================
\section{Слайд 13. Эмпирический GOOG vs смоделированный эталон ОУ}
% ================================================================

Это прямое визуальное сравнение.

Оранжевая линия --- реальная детрендированная лог-цена GOOG ($\widetilde{X}_{t_i}^k$)
с 2023 по 2026 год.
Синяя линия --- одна симуляция ОУ с параметрами, оценёнными на GOOG.

Траектории не идентичны --- это две разные реализации процессов
со схожими статистическими свойствами.
Сравнение носит иллюстративный и учебный характер, это не статистический тест.
Но визуальное сходство поддерживает интерпретацию ОУ.

% ================================================================
\section{Слайд 14. Критерий оценки модели}
% ================================================================

Этот слайд непосредственно отвечает на методологическое замечание профессора Фаткуллина.

Вопрос был: что означает «совместимо»?
Нельзя просто нарисовать две похожие картинки и сказать, что модель подходит.
Нужен явный критерий.

Использованный подход состоит из трёх частей.

\textbf{Первое:} подгоняем параметры по одной простой статистике ---
регрессия AR(1) на детрендированных уровнях,
дающая ($\hat\alpha$) и ($\hat\phi$).

\textbf{Второе:} проверяем другие следствия подогнанной модели,
которые \emph{не использовались} при подгонке:
\begin{itemize}
  \item Проверка первая: соответствует ли теоретическое убывание ($\hat\phi^\ell$)
        эмпирической АКФ на лагах от 2 до 20?
  \item Проверка вторая: предсказывает ли условное среднее модели ---
        ($\hat m + \hat\phi^h(\widetilde{X}_t - \hat m)$) ---
        эмпирическое условное среднее при $h = 5$ дней?
\end{itemize}

\textbf{Третье:} фиксируем видимый изъян ---
нестационарную волатильность.

Это не полный тест на переходное ядро, как также предлагал профессор.
Но это систематическая диагностика, дающая явный смысл
слову «совместимо».

% ================================================================
\section{Слайд 15. Оценка модели: количественная таблица}
% ================================================================

Таблица показывает все метрики вместе.

\textbf{Коэффициент детерминации тренда} ($R^2$) высок для всех трёх акций:
0.883 для NVDA, 0.822 для META, 0.866 для GOOG.
Эталон геометрического броуновского движения применим к сырым уровням лог-цен.

\textbf{Коэффициент детерминации AR(1) для ОУ} составляет примерно 0.98
для всех трёх акций: одношаговая подгонка очень точная.

\textbf{Среднеквадратичная ошибка АКФ} (ACF RMSE)
мала, особенно для GOOG --- 0.007:
теоретическая кривая убывания ($\hat\phi^\ell$) хорошо совпадает
с эмпирической АКФ.
Для NVDA отклонение больше --- 0.049:
этот ряд менее точно укладывается в эталон ОУ.

\textbf{СКО переходного условного среднего}
составляет примерно от 0.005 до 0.007:
предсказание условного среднего точное.

\textbf{Отношение скользящей волатильности} превышает 3--7 раз для всех акций:
именно здесь СДУ с постоянной сигмой явно не справляется.
META показывает наихудшее отношение --- 6.8 раза:
у неё был особенно турбулентный период в 2023 году.

% ================================================================
\section{Слайд 16. Основные выводы}
% ================================================================

Ответ взвешенный, не абсолютный.

\textbf{Геометрическое броуновское движение (GBM)} подходит для сырых лог-цен ---
они доминированы трендом, растут со временем,
и картина мультипликативного шума хорошо согласуется.

\textbf{Орнштейн--Уленбек (OU)} подходит для детрендированных лог-цен ---
они показывают персистентную, медленно убывающую автокорреляцию
с периодом полуравновесия около трёх месяцев,
и обе проверки --- убывание АКФ и условное среднее --- проходят.

\textbf{Диффузия с постоянной сигмой} полезна, но ограничена ---
язык Ито естественно описывает дневные приращения,
но волатильность явно нестационарна.

Ни одно низкоразмерное стохастическое дифференциальное уравнение
не улавливает все свойства одновременно.

% ================================================================
\section{Слайд 17. Ограничения и естественные расширения}
% ================================================================

Основные ограничения следующие.

Линейное детрендирование может быть неоптимальным способом удаления тренда;
факторное детрендирование могло бы дать другие результаты.

Проверка условного среднего --- это упрощённая диагностика,
а не полный анализ переходного ядра, упомянутый профессором Фаткуллиным.

СДУ с постоянной волатильностью полностью упускает кластеризацию волатильности.

Наиболее естественный следующий математический шаг ---
\textbf{реализованная квадратичная вариация}:
для диффузии со сносом $u$ и коэффициентом диффузии $v$
квадратичная вариация ($\langle X\rangle_t$)
сходится к ($\int_0^t v^2\,ds$).
На практике сумма квадратов дневных доходностей даёт прямую эмпирическую оценку
интегрированной дисперсии.
Это была бы более строгая диагностика диффузии, чем скользящая волатильность.

Ещё одно естественное направление --- \textbf{уравнение Фоккера--Планка}
для процесса ОУ,
которое задаёт эволюцию плотности вероятности ---
это позволило бы провести более формальный анализ переходного ядра,
как предложил профессор.

% ================================================================
\section{Слайд 18. Заключение}
% ================================================================

Заключение намеренно взвешенное, не абсолютное.

Исследовательский вопрос был: могут ли простые стохастические дифференциальные уравнения
описывать временные ряды фондового рынка?
Ответ: да, в значимой мере,
но лишь приближённо и только для правильных преобразований данных.

Методологический урок состоит в том, что «совместимо» --- это не визуальное суждение.
Подгоняете по одним статистикам, проверяете по другим.

Это учебное мини-исследование, и я считаю, что ценность его состоит именно
в том, чтобы сделать связь между теоретическим аппаратом Оксендаля
и реальными рыночными данными явной ---
показать, какие формулы применимы где, и каковы их ограничения.

\end{document}
